%% notes-tex/interaction-fitness-figure-gate/sections/2-mockups.tex
%% SOURCE: the ladder sheet is written by experiments/030-solid-growth-multi/scripts/
%% lower_order_baselines_030.py and the fitness panel by figure2_gate_numbers.py; every
%% number in a caption is a macro of tables/gate_numbers.tex written by the latter.

\section{The mockups\secstatus{todo}}
\label{sec:mockups}

Result files are under \file{experiments/030-solid-growth-multi/results/} unless a path
is given. Each caption gives, per panel, what is plotted, the statistic, the sample size
and the status. The statistic is the Pearson correlation on the pinned test triples of the
arm unless the caption says otherwise; a transformer row is a validation number, the mean
over the last ten epochs synced to the run registry, because no 030 checkpoint has a test
evaluation yet. \regimekey{}

\pfstep{The sheet proposed for the main text} Seven panels in three rows, in the order
the work happens: the setup, the measured data, what a model reads. Panels a and b are
compositions with no data behind them yet; c to e are on file; f and g wait on runs.

\begin{description}
  \item[a] What a trigenic interaction is: the identity of \ref{sec:decision} drawn as the
    three singles, the three doubles and the triple of one Kuzmin query pair, with the
    terms that are training records marked and the one that is not. \fstat{idea}
  \item[b] What a record is on the 030 build and how the two arms split it: a genotype
    with its entries from every screen, the closure pool S3 (singles, closure doubles,
    triples), arm R pinned over triples, arm Q pinned over query pairs with the held-out
    pairs' doubles removed. \fstat{idea}
  \item[c, d] The interaction ladder on arm R and arm Q, Figure~\ref{fig:ladder}. \fstat{file}
  \item[e] The fitness side, Figure~\ref{fig:fitness}. \fstat{file}
  \item[f] The transformer with the lower-order terms as inputs, beside the one that reads
    identity only, three seeds each, test: the regime comparison. \fstat{idea}
  \item[g] Accuracy against $|\tau|$ and the scaling readout of \file{030-scaling-study}.
    \fstat{idea}
\end{description}

\clearpage
\tcfig{\figLadder}{\textbf{The interaction ladder on both arms: what the lower-order
terms explain, and where the transformer sits.} \textbf{a}, arm R, the 010 random split
over triples; \textbf{b}, arm Q, the query-pair-disjoint split. One bar per rung, test
Pearson on \gnRTest{} (a) and \gnQTest{} (b) held-out triples in the arm's light fill;
the last bar is the 030 transformer's validation mean over its last ten synced epochs in
the arm's darker color (a: epochs \gcgtRFirstEpoch{} to \gcgtRLastEpoch{} of seed
\gcgtRSeed{}; b: epochs \gcgtQFirstEpoch{} to \gcgtQLastEpoch{}). Rungs, left to right:
the screen mean (\gloRNZeroscreenmeantest{}, \gloQNZeroscreenmeantest{}), the
zero-parameter digenic-controls null (\gloRNOnemultquerydoubletest{},
\gloQNOnemultquerydoubletest{}), the singles-product null (\gloRNTwomultsinglestest{};
masked on Q, where it falls to the floor), the strain-matched controls
(\gloRNOnesmultstrainmatchedtest{}, \gloQNOnesmultstrainmatchedtest{}), the additive
gene ridge (\gloRBOneadditivegenetest{}, \gloQBOneadditivegenetest{}), ridge over the
lower-order terms (\gloRLOneridgelowerordertest{}, \gloQLOneridgelowerordertest{}),
boosted trees over the terms (\gloRLTwohgblowerordertest{},
\gloQLTwohgblowerordertest{}), ridge over genes and terms
(\gloRLThreeridgegenepluslowerordertest{}, \gloQLThreeridgegenepluslowerordertest{}),
and the transformer (\gcgtRValWindow{}, \gcgtQValWindow{}). Written by
\texttt{lower\_order\_baselines\_030.py}, GilaHyper job 3902. \fstat{file}}{fig:ladder}

\noindent\textbf{What each rung of Figure~\ref{fig:ladder} reads.} The screen mean reads
the training mean of $\tau$ in the triple's screen, Kuzmin 2018 or 2020, and is the
floor every pooled Pearson carries. The digenic-controls null predicts
$-(\varepsilon_{ik} + \varepsilon_{jk})$, which equation~\ref{eq:identity} gives when the
triple's fitness equals its multiplicative expectation $f_{ij} f_k$; it has no parameters
and reads two training records. The singles-product null replaces $f_{ijk}$ by
$f_i f_j f_k$, equal to the first null minus $\varepsilon_{ij} f_k$, and on arm Q it has
no $f_{ij}$ to read. The strain-matched rung takes each control score from the entry
measured against the triple's own array strain in its own screen, as the source computed
it. The additive gene ridge is the 010 rung, one coefficient per gene. The ridge over
the lower-order terms uses thirteen terms, four products, a screen indicator and a
missing indicator per column; the boosted trees use the same terms with missing values
handled natively, 600 trees, leaf count and learning rate chosen on validation; the last
ridge stacks the gene one-hot beside the terms. Where a term is missing a null falls
back to the screen mean for that row, so every rung is scored on identical rows, and
the performance table carries the covered fraction.

\tcfig{\figFitness}{\textbf{The fitness side is arithmetic.} Pearson on the held-out
triples' own fitness $f_{ijk}$: a ridge over the lower-order terms (test, light fill;
\gloRFOneridgelowerorderfitnesstest{} on arm R over \gnRTest{} triples,
\gloQFOneridgelowerorderfitnesstest{} on arm Q over \gnQTest{}) beside the 030
transformer's fitness head on the same triples (validation, darker; \gfitRCgtValWindow{}
on R over epochs \gfitRCgtFirstEpoch{} to \gfitRCgtLastEpoch{}, \gfitQCgtValWindow{} on
Q over epochs \gfitQCgtFirstEpoch{} to \gfitQCgtLastEpoch{}). Written by
\texttt{figure2\_gate\_numbers.py} from the ladder file and the run registry.
\fstat{file}}{fig:fitness}

\pfstep{The supplement} Three sheets are proposed, none drawn yet. First, how the query
decides the number: the same model family on the 010 build (triples only), the 025 S3
closure pool (merged labels) and the 030 per-entry pool (every entry, a dataset token),
one bar per build per arm, with the gene-identity ladder of \file{025-additive-baselines}
beside it; what is on file is one seed per build and the ladders. \fstat{partial} Second,
the screen-offset floor: per-screen and within-screen Pearson of every rung and the
transformer, the per-screen mean $\tau$ that makes the floor, and the Kuzmin 2018 against
2020 difference in what the arithmetic explains (the controls null reads
\gloRNOnemultquerydoubletestKzTwoZeroOneEight{} on 2018 rows and
\gloRNOnemultquerydoubletestKzTwoZeroTwoZero{} on 2020 rows of arm R test). \fstat{file}
Third, the comparators of the earlier draft on the same data: DCell, DANGO and Yeast9 FBA
refit or rescored on the 030 pinned triples of both arms. \fstat{idea}
